\begin{figure}
\TwoColDIY%
    {.6}{\EquivalenceDef{}
          \SetoidDef{}}%
    {.35}{%
      \SetoidSugar{}
      \SetoidCarrierSugar{}
      \SetoidRelationSugar{}
    }
    \caption{(a) Equivalence relations and setoids as records and
      (b) example desugaring into a GADT and projections}
\figlabel{setoid}
\end{figure}
\paragraph{Setoids in \idris{}.}
We represent equivalence
relations and setoids in \idris{} with \emph{records} in \figref{setoid}a.
\idris{} records are syntactic sugar for a single-constructor
\IdrisKeyword{data} declaration and automatically generated
\emph{field projections}, as in \figref{setoid}b.  \idris{} also
automatically generates the post-fix projections for each
field using a dotted notation, writing
\IdrisBound{b}\IdrisFunction{.equivalence.relation} for the
nested projection.
The annotation \IdrisKeyword{0} preceding the definition of the field $\U$ is a
\emph{quantity} \cite{mcbride:qtt,atkey:qtt} indicating
that the field is not represented at runtime, but may be used in
types. Readers can safely ignore these annotations. There is only a handful of them in
this article.
We maintain them to demonstrate and emphasise that \Frexlib{} does not
require us to retain, for example, runtime representations of types.
If you are reading this article in colour,
our listings include semantic highlighting, designating the semantic class
of each lexeme: \IdrisData{data} constructor, \IdrisType{type}
constructor, \IdrisFunction{defined} function or value, and
\IdrisBound{variable} in a binding/bound occurence.
We define setoid homomorphisms:
\MyCodeListing{
  \PutListing{.49\textwidth}\SetoidHomo
  \PutListing{.39\textwidth}\SetoidArrow
}
The Appendix includes expanded examples for setoids of functions and quotient
setoids.

This technique is affectionately dubbed `setoid hell', since we
need to prove that all our functions are setoid homomorphisms.
Following \citeauthor{hu-carette:agda-ct}~\citeyearpar{hu-carette:agda-ct}, we manage
setoid hell by structuring code categorically, organising results
into homomorphisms between appropriate setoids.
